\chapter{CSO流塑造拓扑：基于MDL的Agent-relative条件结构复杂度理论}
\label{chap:cso-mdl-topology}

\section{问题的提出：从“复杂度迁移”回到可测量的认知结构代价}

前文已经建立了两个相互区别而又彼此关联的基本对象：其一是认知结构本体
\[
C\in\mathcal{C},
\]
本文统一称为认知结构对象（Cognitive Structure Object, CSO）；其二是某一具体智能体对该CSO的内部承载，即Agent-CSO。若智能体记为$A$，则同一个CSO $C$在智能体$A$中的具体实现可以抽象表示为
\[
\operatorname{ACSO}_{A}(C,t)
=
\left(
C,
R_t,
F_t,
\tau_t,
B_t,
S_t
\right),
\]
其中$R_t$表示具体表示形式，$F_t$表示当前主要功能承载位置，$\tau_t$表示有效时间尺度，$B_t$表示物理或逻辑载体，$S_t$表示其激活、潜伏、预测或执行状态。由此可以看出，一个CSO是否改变并不能仅由其内部编码是否变化来判定；同一认知结构完全可能在不失去结构连续性的情况下，由显式表示迁移为长期生成结构，由高层认知过程迁移为局部技能，或者由内部记忆迁移至外部工具。

在早期理论中，这类现象被概括为“复杂度迁移”。例如，一个原本需要大量显式推理的问题，在反复学习后变得自动；一个需要长期内部维持的信息被写入文件或程序；两个功能区域之间反复发生的信息交互，最后形成稳定的直接连接。从直觉上看，可以说“复杂度从一个地方迁移到了另一个地方”。然而严格来说，复杂度本身并不是一种像能量一样能够独立于结构存在和流动的实体。真正发生迁移的是CSO的表示、承载、时间尺度以及功能位置，而所谓复杂度变化，是这种迁移造成的描述成本、运行成本和关系维护成本的重新分布。

因此，本章将“复杂度迁移”重新建立在最小描述长度（Minimum Description Length, MDL）和条件算法复杂度的思想上，并定义一种面向具体智能体的条件CSO复杂度：
\[
\boxed{
L_A(C\mid M_t,B)
}
\]
它刻画的是：在智能体$A$已经拥有某些内部结构$M_t$的条件下，如果要求将CSO $C$承载在载体$B$中，还需要增加多少新的有效描述。该量并不试图计算不可计算的严格Kolmogorov复杂度，而是把Kolmogorov复杂度作为理论极限，把MDL作为工程上可以近似优化的结构准则。

这样，前文所提出的
\[
\text{CSO流塑造拓扑}
\]
可以获得一个更加严格的数学解释：长期反复发生的CSO跨功能迁移，如果能够通过形成稳定连接而减少未来总描述长度、协调成本和实时负担，那么动态关系就具有沉降为静态拓扑的长期优势。拓扑学习因此可以理解为历史CSO流的一种结构压缩。

本章将首先定义Agent-relative条件CSO复杂度，然后建立有效认知成本、工作记忆有限容量、结构迁移、记忆沉降、技能编译、外部化和拓扑学习之间的统一数学关系。后续章节可以在此基础上分别展开三相记忆、工作记忆管理、技能系统和具身运行时，但本章并不依赖这些后续具体实现，而只使用它们的最小抽象形式。

\section{从Kolmogorov复杂度到Agent-relative条件描述长度}

设$U$为固定的通用计算机。对象$x$的Kolmogorov复杂度定义为
\[
K_U(x)
=
\min_{p:U(p)=x}|p|,
\]
其中$p$是能够生成$x$的程序，$|p|$为程序的描述长度。忽略不同通用计算机之间的常数差异后，可以简写为$K(x)$。它所回答的问题是：一个结构最短可以被描述到什么程度。

对于认知系统而言，更重要的通常不是无条件复杂度$K(C)$，而是条件复杂度
\[
K(C\mid M),
\]
即在已经给定结构$M$的条件下，描述$C$还需要增加多少信息。如果一个智能体已经掌握了某类稳定规律，那么新的实例对它而言通常具有很低的条件复杂度。例如一个已经掌握二次方程求解规则的智能体，面对一个新的二次方程时，并不需要重新编码完整求解理论，只需要编码新的系数和少量任务相关状态。因此，即使环境中的本体CSO没有改变，也可能存在
\[
K(C\mid M_{t+1})
<
K(C\mid M_t).
\]

这给出了学习的一个非常重要的结构解释：学习并不必然使世界本身变得更加简单，而是使世界中的重要结构相对于智能体已有内部结构变得更加容易描述。

严格Kolmogorov复杂度存在不可计算性，因此不能直接作为AgentOS运行时变量。为此，设智能体$A$在载体$B$上允许使用的编码模型族为
\[
\mathfrak{Q}_{A,B}.
\]
对于候选编码$q\in\mathfrak{Q}_{A,B}$，设对应解码器为$D_B$。定义智能体$A$在时间$t$相对于已有结构$M_t$、使用载体$B$承载CSO $C$时的条件MDL复杂度为
\[
\boxed{
L_A(C\mid M_t,B)
=
\min_{q\in\mathfrak{Q}_{A,B}}
\left\{
\ell(q)
\;:\;
D_B(q,M_t)=C
\right\},
}
\]
其中$\ell(q)$是使用某个具体前缀编码体系得到的描述长度。

这里的$M_t$不必预先绑定某一具体认知架构。它只表示智能体当前已经能够作为先验使用的结构集合。可以一般地写成
\[
M_t
=
\left(
M_t^{\mathrm{fast}},
M_t^{\mathrm{episodic}},
M_t^{\mathrm{structural}},
M_t^{\mathrm{self}},
M_t^{\mathrm{goal}},
\mathcal{T}_t
\right),
\]
其中$\mathcal{T}_t$是当前功能拓扑。后续AgentOS实现可以分别把这些结构具体化为工作记忆、情景记忆、长期生成结构、自我结构以及目标结构，但本章只使用它们作为条件编码背景。

由于不同智能体的$M_t$不同，一般有
\[
L_A(C\mid M_t^A,B)
\neq
L_B(C\mid M_t^B,B).
\]
因此必须明确：

\[
\boxed{
\text{CSO本体复杂度}
\neq
\text{Agent-CSO条件复杂度}.
}
\]

前者关注结构本身可以被多短地描述，后者关注该结构对于某个具体智能体还剩下多少新增描述负担。

\section{Agent-relative复杂度的基本性质}

条件CSO复杂度具有三个对智能理论十分重要的性质。第一，它是主体相关的。第二，它随生命周期变化。第三，它与载体相关。

对同一个CSO $C$，如果智能体已经形成能够高度解释它的内部结构$M_t$，则
\[
L_A(C\mid M_t,B)
\ll
L_A(C\mid \varnothing,B).
\]
因此学习可以表现为
\[
L_A(C\mid M_{t+1},B)
<
L_A(C\mid M_t,B).
\]
这可以称为\textbf{条件结构压缩}。

另一方面，即使条件结构不变，同一个CSO在不同载体中的成本也可能不同：
\[
L_A(C\mid M_t,B_i)
\neq
L_A(C\mid M_t,B_j).
\]
例如，某一结构在显式工作状态中可能需要完整展开，而在长期生成结构中只需要一个参数化生成器；在外部程序中则可能只需要一个稳定调用句柄。因此，智能体的学习不仅表现为$M_t$本身发生变化，也表现为对载体$B$的重新选择。

进一步，Agent-CSO的变化可以定义为
\[
z_t(C)
=
\left(
R_t,F_t,\tau_t,B_t,S_t
\right),
\]
从而一次认知结构迁移写作
\[
\boxed{
\mathcal{M}_C:
z_t(C)
\longrightarrow
z_{t+1}(C).
}
\]
这种迁移可以仅改变表示$R$，也可以改变功能位置$F$、时间尺度$\tau$、载体$B$或者激活状态$S$。因此，后文所谓CSO迁移并不意味着本体对象$C$必须变化，而是指其Agent-relative承载方式发生变化。

\section{任务相关失真与最小充分CSO表示}

如果要求智能体无损地保存现实中的全部结构，则任何有限系统都会迅速失去可行性。智能真正需要保存的不是现实的逐比特复制，而是对于当前和未来任务足够的结构。为此必须把MDL与任务相关失真结合。

设$C$是真实或目标CSO，$\hat C$是智能体实际保存的近似结构。定义Agent-relative任务失真函数
\[
d_A
\left(
C,\hat C
\mid
G_t,\Psi_t
\right),
\]
其中$G_t$表示当前目标约束，$\Psi_t$可以一般地理解为主体相关慢结构。若不希望在本章引入具体的自我模型，也可以把二者合并写成上下文约束$\Gamma_t$：
\[
d_A(C,\hat C\mid\Gamma_t).
\]

给定允许误差$\varepsilon$，定义有损条件CSO复杂度
\[
\boxed{
L_A^\varepsilon(C\mid M_t,B)
=
\min_{\hat C,q}
\ell(q)
}
\]
满足
\[
D_B(q,M_t)=\hat C,
\]
并且
\[
d_A(C,\hat C\mid\Gamma_t)
\leq
\varepsilon.
\]

因此智能压缩真正解决的问题不是
\[
\min L,
\]
而是
\[
\boxed{
\min L
\qquad
\text{s.t.}
\qquad
d_A(C,\hat C)\leq\varepsilon.
}
\]

可以进一步定义任务充分性函数
\[
\operatorname{Suff}_A(\hat C,G_t),
\]
要求
\[
\operatorname{Suff}_A(\hat C,G_t)
\geq
\sigma_{\min}.
\]
于是可以把Agent认知压缩理解为寻找
\begin{statementbox}
Minimum Sufficient Representation.
\end{statementbox}

这一点十分关键。智能并不追求最短描述本身，而追求在任务、现实和安全约束下的最短充分描述。若仅仅追求
\[
\min L(M),
\]
最优解可能退化为几乎不包含任何有用信息的常数模型。因此后续所有MDL优化都必须与残差、任务失真和现实反馈共同考虑。

\section{有效条件CSO成本：从描述长度到真实认知负担}

单纯的描述长度仍然无法刻画一个实时智能体所面对的全部代价。一个极短的生成程序如果每次需要极长计算时间才能展开，对实时控制而言可能远劣于一个略长但能够常数时间访问的结构。因此，本章进一步定义\textbf{有效条件CSO成本}：
\[
\boxed{
\mathscr{C}_A(C,t;B)
=
L_A^\varepsilon(C\mid M_t,B)
+
\lambda_T T_B(C)
+
\lambda_W W_B(C)
+
\lambda_R R_B(C)
+
\lambda_E E_B(C).
}
\]

其中$T_B(C)$表示从载体$B$访问、解码或执行该结构需要的时间成本，$W_B(C)$表示实时工作状态占用，$R_B(C)$表示检索、通信、恢复和可靠性相关成本，$E_B(C)$表示能量、经济资源或外部动作等附加成本；$\lambda_T,\lambda_W,\lambda_R,\lambda_E$是由当前系统状态和任务条件决定的权重。

因此必须始终区分两个量：
\[
\boxed{
L_A(C\mid M_t,B)
}
\]
是结构描述长度，而
\[
\boxed{
\mathscr{C}_A(C,t;B)
}
\]
是智能体实际承担该CSO时的综合有效成本。

前者连接Kolmogorov复杂度和MDL，后者才直接参与Agent的在线调度、迁移和载体选择。

可以把不同载体上的有效成本看成一个非物理的认知势函数：
\[
V_A(C,B,t)
=
\mathscr{C}_A(C,t;B).
\]
若存在两个载体$B_i$和$B_j$，满足
\[
V_A(C,B_j,t)
+
C_{\mathrm{mig}}(B_i,B_j)
<
V_A(C,B_i,t),
\]
则从局部成本意义上存在将CSO从$B_i$迁移至$B_j$的动力。

这里的$V_A$只是优化意义上的成本势，而不是物理能量。它的作用仅在于统一描述不同承载状态之间的选择关系。

\section{工作记忆有限容量的MDL形式}

设智能体当前需要显式维护的Agent-CSO集合为
\[
H_t
=
\{C_1,C_2,\ldots,C_n\},
\]
并且这些结构之间存在必须同时维持的关系集合或关系超图
\[
\mathcal{R}_t.
\]
则工作状态真正需要维护的联合描述不是简单的
\[
\sum_i L(C_i),
\]
而是
\[
\boxed{
L_h(H_t)
=
L_A
\left(
C_1,\ldots,C_n,\mathcal{R}_t
\mid
M_t
\right).
}
\]

该表达式强调，认知负荷的主要来源之一是关系结构而非对象数量本身。同样数量的对象，如果彼此独立，其联合描述长度可能相对较低；如果它们之间存在大量条件依赖、冲突、目标约束和动态关系，则
\[
L(\mathcal{R}_t\mid C_1,\ldots,C_n)
\]
可以迅速增长。

设工作记忆在给定最大访问时延$T_{\max}$下能够稳定区分并实时操作的有效内部状态数为
\[
|\mathcal{S}_h(T_{\max})|.
\]
定义其实时有效描述容量
\[
\boxed{
B_h(T_{\max})
=
\log_2
|\mathcal{S}_h(T_{\max})|.
}
\]

则工作记忆的可行条件为
\[
\boxed{
L_h(H_t)
\leq
B_h(T_{\max}).
}
\]

\textbf{定理（有限工作结构编码定理）}
若一个工作记忆系统在时间窗口$T_{\max}$内只能稳定区分有限个内部状态
\[
|\mathcal{S}_h(T_{\max})|<\infty,
\]
且要求从工作状态无歧义地区分所有联合认知状态$(H_t,\mathcal{R}_t)$，则任何无损编码均必须满足
\[
L_h(H_t)
\leq
\log_2|\mathcal{S}_h(T_{\max})|.
\]

\textbf{证明：}
若系统需要区分$N$个不同的联合认知状态，则任何唯一可译码的表示至少需要能够区分$N$种编码状态。若
\[
N>|\mathcal{S}_h(T_{\max})|,
\]
由鸽巢原理，至少有两个不同的联合认知状态被映射到同一个可访问工作状态，因此不能无歧义恢复。由于二进制编码区分$N$种状态至少需要$\log_2N$位，因此可实时稳定维持的联合描述长度上界不超过
\[
\log_2|\mathcal{S}_h(T_{\max})|.
\]
证毕。

该定理的重要意义并不在于给出一个固定数字，而在于说明：只要实时状态空间和允许响应时间有限，就必然存在显式认知承载上限。因此，一个长期智能系统不能依靠不断扩大在线显式状态来解决所有问题，而必须发展出压缩、时间尺度分离、迁移、技能化和外部化机制。

\section{三种基本记忆承载作为MDL最优解的不同区域}

为了避免过度依赖后续具体AgentOS模块，本章只引入三种一般性的认知承载形式：在线工作结构$H$、时间轨迹结构$E$以及长期生成结构$\Theta$。后续章节可以进一步把它们具体化为三相记忆系统，但本章只关心其MDL性质。

在线工作结构$H$强调低访问时延。其目标函数可以表示为
\[
\boxed{
\mathcal{J}_{H}(C)
=
L_H(C\mid M_t)
+
\alpha_H T_H(C)
+
\beta_H W_H(C),
}
\]
其中实时性权重$\alpha_H$较大，因此$H$可以容忍一定表示冗余，以换取快速访问。

时间轨迹结构$E$承担的是无法立即被一般规律解释、但未来仍可能具有价值的时间残差。给定候选长期生成结构$\Theta$，经验序列的总编码可以写成
\[
\boxed{
L(\Theta)
+
L(E\mid\Theta).
}
\]
其中$L(E\mid\Theta)$表示当前生成结构仍无法解释的具体情景残差。因此可以把$E$理解为
\[
\boxed{
\text{结构模型}+\text{时间残差}
}
\]
中的后者。

长期生成结构$\Theta$的任务则是寻找能够用更短描述解释大量重复经验的生成器：
\[
\boxed{
\Theta^*
=
\arg\min_{\Theta}
\left[
L(\Theta)
+
L(E\mid\Theta)
\right].
}
\]
若进一步考虑主体连续性、安全和长期稳定性，可以写成
\[
\Theta^*
=
\arg\min_{\Theta}
\left[
L(\Theta)
+
L(E\mid\Theta)
+
\lambda_S\mathcal{R}_{S}(\Theta)
\right],
\]
其中$\mathcal{R}_S$代表结构更新带来的稳定性代价。

因此三种基本承载形式可以用如下方式统一：
\[
\boxed{
H
=
OnlineExpansion,
}
\]
\[
\boxed{
E
=
TemporalResidualCompression,
}
\]
\[
\boxed{
\Theta
=
ReusableGenerativeCompression.
}
\]

于是从在线工作结构到经验，再到长期生成结构的过程
\[
H\rightarrow E\rightarrow\Theta
\]
可以解释为
\[
\boxed{
\text{实时展开}
\rightarrow
\text{轨迹残差压缩}
\rightarrow
\text{可复用生成压缩}.
}
\]

反向过程
\[
\Theta\rightarrow H
\]
则不是简单复制，而是条件解码：
\[
\hat C_t
=
D_H(q,\Theta,G_t),
\]
即长期压缩结构根据当前任务重新生成适合当前工作状态的显式CSO。

\section{经验向生成结构沉降的MDL判据}

设智能体当前长期模型为$\Theta_0$，候选更新模型为$\Theta_1$，近期经验为$E$. 定义使用旧结构编码经验的总代价
\[
J_0
=
L(\Theta_0)
+
L(E\mid\Theta_0),
\]
以及更新后的总代价
\[
J_1
=
L(\Theta_1)
+
L(E\mid\Theta_1).
\]

定义生成结构更新获得的描述收益
\[
\boxed{
\Delta_\Theta
=
J_0-J_1.
}
\]

然而长期结构改变并不是免费的。设更新成本为$C_{\mathrm{update}}$，已有能力重新校准成本为$C_{\mathrm{recal}}$，长期稳定性代价为$C_{\mathrm{stability}}$。只有满足
\[
\boxed{
\Delta_\Theta
>
C_{\mathrm{update}}
+
C_{\mathrm{recal}}
+
C_{\mathrm{stability}}
}
\]
时，才具有将$\Theta_0$更新为$\Theta_1$的长期MDL优势。

由此可以看到，稳定经验并不会因为出现一次就立即进入长期结构。只有当某种经验模式能够在足够长时间内降低总体编码成本时，它才值得被结构化。这一原则也解释了为什么慢结构应该比快速工作状态拥有更低更新频率：改变越深的结构，其一次性迁移和重新校准成本越高。

\section{工作记忆管理作为在线最小充分描述问题}

设当前完整候选结构为$\mathcal{C}_t$，由于有限容量约束，不可能全部进入在线工作集合$H_t$. 因此智能体需要寻找一个压缩工作表示$\hat H_t$，使
\[
\boxed{
\hat H_t^*
=
\arg\min_{\hat H_t}
L_h(\hat H_t\mid M_t)
}
\]
满足
\[
L_h(\hat H_t)\leq B_h
\]
以及
\[
d_A(H_t,\hat H_t)\leq\varepsilon.
\]

因此，任何后续具体的工作记忆管理机制，无论采用注意、分解、摘要、折叠、暂停、恢复还是局部外化，其数学目标都可以统一为：

\begin{statementbox}
在有限在线描述预算内，维持最小任务充分结构。
\end{statementbox}

这说明工作记忆调度并不是一种附加优化，而是有限认知系统的必需过程。由于$B_h$有限，当输入和任务结构持续增长时，系统必须在以下操作中至少选择一种：压缩当前表示、把部分结构迁移到较慢载体、把重复操作编译成技能、把部分结构迁移到外部载体，或者改变功能拓扑使未来协调成本下降。

因此，工作记忆有限性是后续多种学习机制共同的结构压力来源。

\section{技能编译作为重复在线描述的摊销压缩}

设某类任务$D$被重复执行$n$次。未形成技能时，每次都需要由在线认知过程显式产生控制结构$\pi_i$。若平均每次显式执行的有效成本为$c_e$，则累计成本为
\[
J_e(n)
=
nc_e.
\]

现在允许建立一个技能结构$S$。其一次性编译、校准和保存成本记为
\[
C_S,
\]
形成技能之后，每次调用的平均成本为$c_s$，其中假设
\[
c_s<c_e.
\]
则总成本为
\[
J_s(n)
=
C_S
+
nc_s.
\]

当
\[
C_S+nc_s<nc_e
\]
时，技能化优于持续显式执行，因此有
\[
\boxed{
n
>
\frac{C_S}{c_e-c_s}.
}
\]

\textbf{定理（重复结构技能编译阈值）}
若某类认知—行动结构重复出现，建立技能的一次性成本有限，并且技能化后单次调用成本严格低于显式执行成本，则必然存在有限重复次数$N_S$，使得当$n>N_S$时，技能化的生命周期累计成本严格低于持续显式推理。

该定理说明，技能编译可以理解为典型的摊销压缩。一个原本必须在工作状态中反复展开的Agent-CSO，被转换成一个更适合快速执行的程序化Agent-CSO：
\[
\boxed{
\operatorname{ACSO}^{\mathrm{explicit}}
\longrightarrow
\operatorname{ACSO}^{\mathrm{procedural}}.
}
\]

因此技能形成并不是增加一个与认知理论无关的新模块，而是CSO改变自身承载形式。其本质效果是使
\[
W_H(C)\downarrow,
\qquad
T(C)\downarrow,
\]
并以一次性的结构建立成本换取未来大量实时计算成本的下降。

\section{Offloading作为跨边界载体优化}

设CSO $C$既可以内部承载，也可以写入外部载体$B_e$。内部承载有效成本记为
\[
\mathscr C_{\mathrm{int}}(C).
\]
外部承载成本包括写入、编码、检索、通信、时延和风险：
\[
\boxed{
\mathscr C_{\mathrm{ext}}(C)
=
L_{\mathrm{ext}}(C)
+
C_{\mathrm{write}}
+
p_{\mathrm{reuse}}C_{\mathrm{retrieve}}
+
C_{\mathrm{latency}}
+
C_{\mathrm{risk}}.
}
\]

若
\[
\mathscr C_{\mathrm{ext}}(C)
<
\mathscr C_{\mathrm{int}}(C),
\]
并且满足
\[
Reliability
\geq r_{\min},
\]
\[
Latency
\leq T_{\max},
\]
\[
Security
\geq s_{\min},
\]
\[
Recoverability
\geq q_{\min},
\]
则将$C$迁移到外部载体在长期成本意义上更优。

因此外部化可以定义为
\[
\boxed{
Offloading
=
ConstrainedCrossBoundaryCarrierOptimization.
}
\]

这里的External并不预设具体实现。它可以是文件、数据库、程序、工具、物理环境或者其他能够被稳定重新访问的外部结构。重要的不是物理位置，而是：Agent是否保留了可靠索引、恢复路径和功能关联。如果一个外部结构无法被重新定位或重新解释，则它不能构成有效的长期外部认知载体。

\section{CSO迁移的统一变分定义}

现在可以统一上述多种现象。设CSO $C$在时间$t$的承载状态为
\[
z_t(C)
=
(R_t,F_t,\tau_t,B_t,S_t).
\]
一次迁移为
\[
z_t(C)
\rightarrow
z_{t+1}(C).
\]
设迁移本身需要付出成本
\[
C_{\mathrm{mig}}
\left(
z_t,z_{t+1}
\right).
\]
定义本次迁移带来的净有效收益
\[
\boxed{
\Delta\mathscr C
=
\mathscr C_A(C,t;z_t)
-
\mathscr C_A(C,t;z_{t+1})
-
C_{\mathrm{mig}}(z_t,z_{t+1}).
}
\]

若
\[
\Delta\mathscr C>0
\]
并且所有任务、安全、身份和实时约束仍然成立，则该迁移在当前预测窗口内具有成本优势。

根据$z_t$中发生改变的维度，可以区分：

\[
R_t\rightarrow R_{t+1}
\]
表示形式迁移；

\[
F_t\rightarrow F_{t+1}
\]
功能承载迁移；

\[
\tau_t\rightarrow\tau_{t+1}
\]
时间尺度迁移；

\[
B_t\rightarrow B_{t+1}
\]
载体迁移；

\[
S_t\rightarrow S_{t+1}
\]
激活状态变化。

因此记忆沉降、召回、技能编译、Offloading和后续拓扑学习都可以统一为CSO迁移的不同特例。

\section{功能迁移与动态协调成本}

现在回到本章的核心问题：CSO流为什么会塑造功能拓扑。

设两个功能节点为$F_i$和$F_j$。某类CSO在任务执行中需要反复从$F_i$迁移到$F_j$。若当前拓扑中不存在稳定直接关系，则每次迁移需要进行动态寻址、状态转换、协议选择、表示适配以及冲突处理。记平均动态协调成本为
\[
c_d.
\]

若在$F_i$和$F_j$之间建立稳定功能耦合$e_{ij}$，则需要一次性支付结构建立成本
\[
C_e,
\]
但之后每次相同迁移只需要平均成本
\[
c_e,
\qquad
c_e<c_d.
\]

若该迁移重复$n$次，则旧拓扑累计成本为
\[
J_0(n)
=
nc_d,
\]
形成稳定连接后的累计成本为
\[
J_1(n)
=
C_e+nc_e.
\]

建立稳定边具有优势的条件是
\[
C_e+nc_e<nc_d,
\]
等价于
\[
\boxed{
n
>
\frac{C_e}{c_d-c_e}.
}
\]

这说明，一个稳定拓扑连接并不是因为两个功能“概念上应该连接”而产生，而是因为它们之间存在足够长期、稳定和有价值的结构流，使得预先支付一次结构成本能够降低未来的大量动态协调成本。

\section{拓扑沉降定理}

\textbf{定理（稳定CSO流的拓扑沉降定理）}

设某类CSO在两个功能节点$F_i$与$F_j$之间重复迁移。若：

\begin{enumerate}
    \item 未建立稳定边时，每次动态迁移的平均有效成本为$c_d$；
    \item 建立稳定边后的平均有效成本为$c_e$；
    \item 有
    \[
    \delta=c_d-c_e>0;
    \]
    \item 建立并校准稳定边的一次性总成本为
    \[
    C_T<\infty;
    \]
    \item 该类迁移在未来保持足够稳定，使平均成本差$\delta$具有正的长期期望；
\end{enumerate}

则存在有限阈值
\[
\boxed{
N_T
=
\left\lceil
\frac{C_T}{\delta}
\right\rceil
}
\]
使得当重复次数$n>N_T$时，形成稳定功能连接的累计有效描述成本严格低于持续动态协调。

\textbf{证明：}

未形成稳定连接时，
\[
J_{\mathrm{dynamic}}
=
nc_d.
\]
形成稳定连接后，
\[
J_{\mathrm{topology}}
=
C_T+nc_e.
\]
若
\[
n>
\frac{C_T}{c_d-c_e},
\]
则
\[
n(c_d-c_e)>C_T,
\]
因此
\[
nc_d>C_T+nc_e.
\]
即
\[
J_{\mathrm{dynamic}}
>
J_{\mathrm{topology}}.
\]
证毕。

该定理给出了下式的严格成本解释：
\[
\boxed{
PersistentCSOMigration
\rightarrow
StableFunctionalCoupling.
}
\]

因此可以重新定义
\[
\boxed{
Topology
=
CompressedHistoryOfStableFunctionalInteraction.
}
\]

也就是说，功能拓扑是过去大量重复动态协调被压缩之后留下的慢结构。

\section{拓扑本身也是一种描述代码}

设智能体功能拓扑为
\[
\mathcal{T}
=
(V,E,\mathcal{L}),
\]
其中$V$为功能节点集合，$E$为连接关系，$\mathcal{L}$为稳定闭环集合。拓扑本身必须具有描述成本
\[
\boxed{
L(\mathcal{T}).
}
\]

因此系统不能通过无限增加边和模块来消除所有动态协调成本，因为
\[
|E|\uparrow
\quad\Longrightarrow\quad
L(\mathcal{T})\uparrow,
\]
并且还会增加同步、维护、校准和安全治理成本。

给定生命周期任务数据$\mathcal{D}_{\mathrm{life}}$，可以定义一个基本的拓扑选择问题：
\[
\boxed{
\mathcal{T}^*
=
\arg\min_{\mathcal{T}}
\left[
L(\mathcal{T})
+
L(\mathcal{D}_{\mathrm{life}}\mid\mathcal{T})
+
C_{\mathrm{coord}}(\mathcal{T})
+
C_{\mathrm{risk}}(\mathcal{T})
\right].
}
\]

第一项惩罚过于复杂的组织结构，第二项衡量该拓扑对生命周期经验和任务的解释能力，第三项衡量实时功能协调成本，第四项限制过度连接造成的安全与治理风险。

因此拓扑学习并不是“架构越复杂越好”，而是寻找
\begin{statementbox}
最小足够功能组织结构.
\end{statementbox}

这与MDL中“模型不能太简单，也不能为了拟合数据无限复杂化”的基本原则完全一致。

\section{模块形成作为内部CSO流的压缩}

设一组功能节点
\[
V_M\subseteq V
\]
之间长期发生高度重复的内部交互。若每次都必须显式协调这些关系，则累计描述成本较高。现在将这一组功能关系封装成一个相对稳定模块$M$。模块需要支付结构成本$L(M)$，但能够减少未来大量边界协调和状态描述。

若
\[
\boxed{
L(M)
+
L(\mathrm{Interactions}\mid M)
<
L(\mathrm{Interactions}),
}
\]
则模块化具有描述优势。

因此可以得到前文所提出的命题
\[
\boxed{
Module
=
StableTopologicalCondensate
}
\]
的MDL解释：模块并非智能系统的第一性对象，而是反复发生的内部功能关系在慢时间尺度上被压缩和固化以后形成的拓扑凝聚体。

从这一观点出发，高层认知能力不应因为被观察到就自动对应一个软件模块。只有当某种功能关系长期稳定、重复并且显式物化能够减少未来总体结构和协调成本时，它才有理由从涌现能力沉降为明确的功能结构。

\section{快CSO流与慢拓扑学习}

设$J_{ij}(t)$表示单位时间内从功能$F_i$向$F_j$迁移的CSO流量或者加权结构流。由于单次迁移可能只是偶然事件，拓扑不应直接响应瞬时$J_{ij}(t)$，而应该响应较慢时间窗口中的统计量：
\[
\boxed{
\bar J_{ij}^{(T)}
=
\frac{1}{T}
\int_t^{t+T}
J_{ij}(s)\,ds.
}
\]

拓扑更新可以抽象写为
\[
\boxed{
\mathcal{T}_{t+T}
=
\mathcal{R}
\left(
\mathcal{T}_t,
\bar J^{(T)},
\Delta L,
C_{\mathrm{risk}}
\right),
}
\]
其中$\mathcal{R}$是受约束拓扑重写算子，$\Delta L$表示预期MDL收益。

于是形成一个典型的快慢系统：
\[
\boxed{
FastCSOFlow
\rightarrow
SlowTopologyLearning
}
\]
以及
\[
\boxed{
SlowTopology
\rightarrow
FutureCSOFlow.
}
\]

这与前文多尺度智能动力学中的一般规律完全一致：
\begin{statementbox}
SlowStructureShapesFastDynamics,
\end{statementbox}
\begin{statementbox}
PersistentFastResidualShapesSlowStructure.
\end{statementbox}

在本章语境中，可以进一步改写为
\begin{statementbox}
TopologyShapesCSOFlow,
\end{statementbox}
\begin{statementbox}
PersistentCSOFlowShapesTopology.
\end{statementbox}

这两个式子构成了本章最核心的闭环。

\section{拓扑迟滞与最小必要重构}

如果形成和删除连接没有成本，系统面对短期任务波动时就可能不断重新布线，从而产生拓扑震荡。因此必须考虑拓扑改变的固定成本。设建立连接成本为
\[
C_{\mathrm{on}},
\]
删除连接成本为
\[
C_{\mathrm{off}},
\]
重建表示和重新校准成本为
\[
C_{\mathrm{recal}}.
\]

只有当预期未来收益满足
\[
\Delta J_{\mathrm{future}}
>
C_{\mathrm{on}}
+
C_{\mathrm{recal}}
\]
时，才应建立新的结构关系；反之，只有删除后的预期收益超过
\[
C_{\mathrm{off}}
+
C_{\mathrm{recal}}
\]
时，才值得移除旧结构。

因此系统自然形成
\begin{statementbox}
TopologicalHysteresis.
\end{statementbox}

拓扑迟滞不是附加经验规则，而是具有固定结构更新成本时的自然结果。

进一步，设改变表示、迁移载体、迁移功能、修改拓扑的一次性代价分别为
\[
C_R,
\qquad
C_B,
\qquad
C_F,
\qquad
C_T,
\]
通常应满足
\[
C_R
<
C_B
\lesssim
C_F
<
C_T.
\]

因此在面对新问题时，系统应优先尝试较低代价的结构调整，而只有当低层调整无法获得足够长期收益时，才进入更深层重构。这给出了\textbf{最小必要重构原则}：
\[
\boxed{
a^*
=
\arg\min_a
\left[
C_a
+
J_{\mathrm{future}}(a)
\right].
}
\]

由此可以得到一条一般演化顺序：
\[
\boxed{
RepresentationAdaptation
\rightarrow
CarrierMigration
\rightarrow
FunctionalMigration
\rightarrow
TopologyReorganization.
}
\]

\section{结构连续性与可逆迁移}

此前“复杂度守恒”容易被误解为复杂度总量严格不变。本章用更准确的\textbf{结构连续性原则}替代这一表述。

设能力$\mathcal{K}$在时间$t$仍然被认为存在。则至少应该存在某一载体$B$以及某种解码过程，使
\[
D_B(R_B,M_t)
=
\hat C
\]
并满足
\[
d_A(C,\hat C)
\leq
\varepsilon.
\]
因此，只要能力仍然存在，相关CSO不必保持原始表示，但必须在某处保持足够的可恢复结构。

由此：
\[
\boxed{
StructuralContinuity
\neq
RepresentationIdentity.
}
\]

一个CSO可以经历
\[
H\rightarrow E\rightarrow\Theta\rightarrow Skill\rightarrow External
\]
等多次迁移，而其功能连续性仍然保持。

另一方面，如果环境分布发生变化，原有承载可能不再是最优。设环境从$P_0$变化为$P_1$，原载体$B_1$在$P_0$下最优，但存在$B_2$使
\[
\mathbb{E}_{P_1}
[
\mathscr C_A(C;B_2)
]
+
C_{\mathrm{mig}}
<
\mathbb{E}_{P_1}
[
\mathscr C_A(C;B_1)
].
\]
则最优结构应允许
\[
B_1\rightarrow B_2.
\]

这说明：
\begin{statementbox}
HealthyIntelligenceRequiresReversibleMigration.
\end{statementbox}

已经技能化的结构可以在环境异常时重新显式化；已经形成的长期结构可以在持续残差出现时重新进入经验层；已经形成的拓扑也必须在长期失配时重新开放为可塑结构。

\section{生命周期MDL泛函}

前述所有局部优化可以统一成一个生命周期目标。设时间范围为
\[
t=0,1,\ldots,T,
\]
时刻$t$需要处理的CSO集合为$\mathcal C_t$，每个CSO的承载状态为$z_t(C)$，功能拓扑为$\mathcal T_t$。定义生命周期总代价
\[
\boxed{
\begin{aligned}
\mathcal J_{\mathrm{life}}
=&
\sum_{t=0}^{T}
\gamma^t
\Bigg[
\sum_{C\in\mathcal C_t}
\mathscr C_A
\left(
C,t;z_t(C)
\right)
\\
&\qquad
+
\sum_{C\in\mathcal C_t}
C_{\mathrm{mig}}
\left(
z_t(C),z_{t+1}(C)
\right)
\\
&\qquad
+
C_{\mathrm{coord}}(t)
+
C_{\mathrm{external}}(t)
+
C_{\mathrm{risk}}(t)
\Bigg]
\\
&
+
\lambda_\Theta L(\Theta_T)
+
\lambda_{\mathcal T}L(\mathcal T_T)
+
\lambda_S L(\mathcal S_T).
\end{aligned}
}
\]

其中$\gamma\in(0,1]$表示时间权重，$\mathcal S_T$代表已经形成的技能或其他程序化结构。

整个优化必须满足：
\[
L_h(H_t)
\leq
B_h,
\]
\[
d_A(C_t,\hat C_t)
\leq
\varepsilon_t,
\]
\[
Latency_t
\leq
T_{\max},
\]
\[
Safety_t
\geq
s_{\min},
\]
以及必要时的主体连续性和权限边界约束。

于是智能体的长期学习可以抽象为
\[
\boxed{
\min_{\{
z_t,
\Theta_t,
\mathcal T_t,
\mathcal S_t,
\mathcal X_t
\}}
\mathcal J_{\mathrm{life}}.
}
\]

这一变分形式并不意味着真实智能体必须显式计算该全局最优解。它更适合作为理论目标函数：各局部控制、记忆、技能、外部化和结构学习机制，都可以被理解成对这一不可直接全局求解问题的多时间尺度近似。

\section{多载体CSO统一定理}

\textbf{定理（多载体CSO最优承载统一定理）}

设持续智能体具有有限或可枚举的候选载体集合
\[
\mathcal B,
\]
每一载体$B\in\mathcal B$对CSO $C$具有有限有效成本
\[
\mathscr C_A(C,t;B),
\]
任意两个允许迁移的载体之间具有有限迁移成本
\[
C_{\mathrm{mig}}(B_i,B_j),
\]
工作记忆满足
\[
L_h(H_t)\leq B_h,
\]
并且任务失真、时延、安全和连续性满足给定约束，则在有限预测区间内，工作记忆分配、经验沉积、长期结构压缩、技能编译和外部化都可以表示为同一个受约束载体分配问题：
\[
\boxed{
\min_{\{B(C,t)\}}
\sum_t
\sum_{C\in\mathcal C_t}
\left[
\mathscr C_A
(C,t;B(C,t))
+
C_{\mathrm{mig}}
\right].
}
\]

若进一步允许功能拓扑$\mathcal T$作为慢变量变化，则拓扑学习同样可以纳入
\[
\boxed{
\min_{\{B(C,t),\mathcal T_t\}}
\mathcal J_{\mathrm{life}}.
}
\]

\textbf{解释：}
工作记忆对应选择低时延、有限容量的在线载体；长期记忆对应选择更慢但更高压缩率的生成载体；技能编译对应将重复显式结构迁移至低调用成本的程序化载体；Offloading对应将结构迁移到外部可恢复载体；拓扑学习则改变不同功能和载体之间未来迁移的成本矩阵。因此它们虽然在工程实现上表现为不同机制，但在数学上属于同一个“CSO在多载体、多功能、多时间尺度空间中的受约束长期承载优化问题”。

\section{从复杂度迁移到条件描述成本重分布}

至此，可以对早期“复杂度迁移”概念作严格修正。假设某CSO从
\[
z_i
\]
迁移到
\[
z_j.
\]
真正发生的是Agent-CSO承载状态变化：
\[
\boxed{
CSORemainsStructurallyRelated,
\qquad
RepresentationAndCarrierChange.
}
\]

由于
\[
L_A(C\mid M_t,z_i)
\neq
L_A(C\mid M_t,z_j),
\]
并且
\[
\mathscr C_A(C,t;z_i)
\neq
\mathscr C_A(C,t;z_j),
\]
从系统层面观察就表现为“复杂度从一个地方转移到了另一个地方”。

因此，更严格的定义应为：
\[
\boxed{
ComplexityMigration
=
ConditionalDescriptionCostRedistribution
InducedByCSOMigration.
}
\]

也就是说，复杂度不是迁移的本体，而是认知结构迁移以后产生的测量结果。真正的本体对象仍然是CSO及其承载关系。

这一区分十分重要，因为它防止把MDL成本误认为某种物理守恒量，同时保留“复杂度迁移”作为一个非常有用的宏观工程描述。

\section{学习、熟练与结构资本的统一定义}

在本章框架下，可以给学习一个比参数更新更一般的定义。设未来重要CSO服从任务分布$P_{\mathrm{future}}$，则若结构更新后满足
\[
\boxed{
\mathbb E_{C\sim P_{\mathrm{future}}}
[
\mathscr C_{A_{t+1}}(C)
]
<
\mathbb E_{C\sim P_{\mathrm{future}}}
[
\mathscr C_{A_t}(C)
],
}
\]
同时任务正确性、现实一致性和安全约束未下降，则可以认为发生了有效学习。

因此：
\[
\boxed{
Learning
=
TransformationOfAgentStructure
ThatReducesExpectedFutureConditionalCSOCost.
}
\]

熟练则可以进一步定义为某一任务分布$D$上的有效成本下降。设
\[
\operatorname{Prof}_A(D)
=
-
\mathbb E_{C\sim D}
[
\mathscr C_A(C)
],
\]
则熟练过程表现为
\[
\operatorname{Prof}_A(D)\uparrow,
\]
尤其伴随
\[
W_H(C)\downarrow,
\qquad
T(C)\downarrow.
\]

因此成熟智能的一个基本特征不是熟悉任务上进行越来越多显式计算，而是：
\begin{statementbox}
越来越多历史动态计算被转换成可复用结构资本.
\end{statementbox}

这些结构资本可以存在于长期生成结构、技能、外部工具以及功能拓扑之中。

从MDL角度看，所谓Topological Capital就是通过一次性支付
\[
L(\mathcal T)
\]
来减少未来反复发生的
\[
L(\mathrm{Coordination}\mid\mathcal T).
\]
因此，拓扑本身可以被理解为一种对历史CSO流的长期编码。

\section{为什么有限工作记忆会推动多尺度智能结构出现}

本章最后可以回答一个更基本的问题：为什么一个智能系统需要发展出不同时间尺度的记忆、技能、外部化以及功能拓扑，而不是永远依靠一个巨大的在线推理空间。

设实时工作结构满足
\[
L_h(H_t)\leq B_h.
\]
如果$B_h$无限大、所有状态访问时间为零、计算资源无限、存储免费、外部世界完全稳定，那么许多结构压缩和迁移确实不再必要。然而实际智能系统满足
\[
\boxed{
FiniteWorkingMemory
+
FiniteTime
+
FiniteEnergy
+
OpenWorld.
}
\]

因此反复出现的结构如果永远停留在高成本在线状态，将持续消耗有限$B_h$。为了维持更大的任务范围，系统必须把不同类型的结构逐渐分配到更适当的承载方式：

\[
\boxed{
\text{短时立即需要}
\rightarrow
\text{在线展开},
}
\]

\[
\boxed{
\text{尚未完全理解但值得保留}
\rightarrow
\text{经验轨迹},
}
\]

\[
\boxed{
\text{跨情境重复规律}
\rightarrow
\text{生成结构},
}
\]

\[
\boxed{
\text{重复行动模式}
\rightarrow
\text{技能},
}
\]

\[
\boxed{
\text{内部长期维持不经济}
\rightarrow
\text{外部载体},
}
\]

\[
\boxed{
\text{长期重复功能迁移}
\rightarrow
\text{拓扑沉降}.
}
\]

因此可以形成一条统一的因果链：
\[
\boxed{
\begin{aligned}
&\text{有限实时承载}
\\
&\Longrightarrow
\text{结构压缩}
\\
&\Longrightarrow
\text{时间尺度分离}
\\
&\Longrightarrow
\text{CSO迁移}
\\
&\Longrightarrow
\text{经验、生成结构、技能与外部化}
\\
&\Longrightarrow
\text{长期稳定功能流}
\\
&\Longrightarrow
\text{拓扑沉降}.
\end{aligned}
}
\]

这说明三相记忆、技能形成、Offloading和拓扑学习虽然在表面上属于完全不同的认知功能，实际上可能源于同一个更加基础的约束：有限智能体必须不断寻找更低的未来条件描述成本。

\section{本章结论：拓扑是历史CSO流的压缩记忆}

本章建立了一套基于MDL的Agent-relative条件CSO复杂度理论。其最基础的对象不是“复杂度实体”，而是CSO及其在具体智能体中的承载状态：
\[
\operatorname{ACSO}_{A}(C,t)
=
(C,R,F,\tau,B,S).
\]

一个CSO对于某个Agent的条件结构复杂度定义为
\[
L_A(C\mid M_t,B),
\]
它近似对应严格条件Kolmogorov复杂度在给定实现模型族中的可计算上界。考虑实际实时系统以后，又进一步定义有效认知成本
\[
\mathscr C_A(C,t;B)
=
L_A^\varepsilon
+
\lambda_TT
+
\lambda_WW
+
\lambda_RR
+
\lambda_EE.
\]

由此，学习不再只被定义为参数变化，而可以定义为：
\begin{statementbox}
使未来重要CSO相对于Agent已有结构变得更短、更快、更稳定地可恢复和可使用.
\end{statementbox}

有限工作记忆条件
\[
L_h(H_t)\leq B_h
\]
使系统不可能让所有结构长期停留在显式在线状态，由此产生时间尺度分离和载体分工。在线结构、经验轨迹和长期生成结构可以分别理解为Online Expansion、Temporal Residual Compression和Reusable Generative Compression；重复显式认知在长期累计成本达到阈值后具有编译成技能的优势；内部结构在外部载体更经济且满足可靠性约束时具有Offloading优势；重复跨功能CSO流则在达到结构盈亏平衡点以后具有沉降为稳定拓扑的长期优势。

因此，本章最重要的结论不是“智能追求最短描述”，而是：

\[
\boxed{
\textbf{
有限智能体在任务充分、现实一致、安全和实时性约束下，持续寻找对未来重要认知结构更低成本的承载方式。
}
}
\]

在这一意义上，
\begin{statementbox}
记忆是结构的时间沉降，
\end{statementbox}
\begin{statementbox}
技能是结构的程序化沉降，
\end{statementbox}
\begin{statementbox}
Offloading是结构的跨边界沉降，
\end{statementbox}
而
\begin{statementbox}
拓扑学习是反复CSO功能迁移的结构沉降。
\end{statementbox}

于是本章最终得到两个互为因果的基本关系：
\begin{statementbox}
TopologyShapesFutureCSOFlow,
\end{statementbox}
\begin{statementbox}
PersistentCSOFlowShapesTopology.
\end{statementbox}

它们共同构成智能结构发育最基本的快慢动力闭环。短时间尺度上，认知结构在既有功能拓扑中流动；长时间尺度上，这些流动留下统计痕迹，并在足够稳定、足够重复且能够降低生命周期总成本时沉降为新的结构关系；新的结构关系又重新改变未来CSO的最短路径、条件描述长度和可实现能力空间。

因此，可以把本章的核心命题最终写为：

\[
\boxed{
\textbf{
功能拓扑不是预先给定的永久容器，而是历史认知结构流在更慢时间尺度上的压缩记忆；而学习则是认知结构不断寻找更优表示、时间尺度、功能位置与载体，并在长期重复中逐渐改变承载这些结构的功能拓扑。
}
}
\]
